% !PRINTED 294
\weberpage{294}
%<<CONT
\[
\begin{aligned}
A_\alpha(\varphi,\psi) &\quad -1,\\
A_\varepsilon(\varphi,\psi) &\quad +1,\\
A'(\varphi,\psi) &\quad -1,\\
E'(\varphi,\psi) &\quad +1.
\end{aligned}
\]

Der doppelte Werth der Charakteristik ist dann, wenn wir der Einfachheit halber die Bezeichnung $(\varphi,\psi)$ weglassen (§.\,95)
\weq{3}{-A_\alpha + A_\varepsilon - A' + E' = 2k;}
dazu die Gleichung (2) addirt, ergiebt
\weq{4}{k = A_\varepsilon,}
d.\,h. \gesperrt{die Charakteristik ist gleich der Anzahl der im Inneren von $f$ gelegenen Schnittpunkte von $\varphi, \psi$.}

Damit haben wir also ein Mittel gefunden, um die Anzahl der Wurzeln der Gleichung
\[
F(z) = 0,
\]
deren geometrische Bilder im Inneren einer beliebig gegebenen Curve liegen, aus der Charakteristik zu bestimmen. Die willkürlich hinzugefügten äusseren Verbindungen spielen hierbei gar keine Rolle mehr und können vollständig unterdrückt werden.

\begin{center}
§. 97.

Bestimmung der Charakteristik.
\end{center}

Um den im Vorigen bewiesenen Satz anwendbar zu machen, müssen wir noch zeigen, wie wir die Charakteristiken wirklich bestimmen können. Denken wir uns zu diesem Zweck die Curve $f$ so in Theilstrecken getheilt, dass in einer von ihnen nur ein Schnittpunkt mit der Curve $\varphi$ oder ein Schnittpunkt mit der Curve $\psi$ liegt, und bezeichnen wir mit $a$ (Fig. 13) den Anfang, mit $b$ das Ende einer solchen Strecke, mit $\xi$ den in ihr liegenden Schnittpunkt von $f$ mit $\varphi$, so wird der Punkt $\xi$ ein Punkt $E(f;\varphi,\psi)$ sein, wenn das Product $\varphi\psi$ in $a$ positiv ist, und ein Punkt $A(f;\varphi,\psi)$, wenn $\varphi\psi$ in $a$ negativ ist; und damit ist die

\weberfigureimage{weber-b1-fig13.png}{0.42}
%CONT>>
